% PAGES: 193-194
% Continues directly after sec06_c2_4.tex (folio 176), which ends mid-section
% 6.3 with "...or by using a linear solver based on the LU decomposition of
% the matrix A." and explicitly leaves no environment open. The paragraph
% below is a fresh sentence, not a continuation of an open environment.

In this section, we show that, for tridiagonal symmetric positive definite matrices,
the LU solver is more efficient (i.e., has a smaller operation count) than the
Cholesky solver.

\medskip
\noindent\textbf{Cholesky linear solvers for tridiagonal spd matrices.}\\
Let $A$ be an $n \times n$ tridiagonal symmetric positive definite matrix. The linear
system $Ax = b$ can be solved using the Cholesky decomposition based solver $x =
\text{linear\_solve\_cholesky}(A,b)$ from Table 6.2. For an efficient implementation,
the Cholesky decomposition, forward substitution, and backward substitution routines
from Table 6.2 are replaced by the specialized routines for tridiagonal matrices as
seen below.

We begin by showing that the Cholesky factor of a tridiagonal symmetric positive
definite matrix is an upper triangular bidiagonal matrix by analyzing how the Cholesky
decomposition simplifies for tridiagonal matrices.\footnote{More generally, the
Cholesky factor of a banded matrix of band $m$ is also banded of band $m$; see an
exercise from the end of this chapter and Stefanica [38] for details.}

Since the symmetric positive definite matrix $A$ is tridiagonal,\footnote{Note that a
sufficient, although not necessary, condition for the tridiagonal symmetric matrix $A$
to be symmetric positive definite is $|A(i,i)| > |A(i,i-1)| + |A(i,i+1)|$ for all $i =
2:(n-1)$, $|A(1,1)| > |A(1,2)|$, and $|A(n,n)| > |A(n,n-1)|$, since in this case $A$ is
strictly diagonally dominant, and any strictly diagonally dominant symmetric matrix is
symmetric positive definite; see Theorem 5.7.} note that
\begin{align}
A(j,k) &= 0, \qquad \forall\, 1 \le j,k \le n \text{ with } |j-k| \ge 2; \\
A(i,i+1) &= A(i+1,i), \qquad \forall\, i = 1:(n-1). \notag
\end{align}

Let $U$ be the Cholesky factor of $A$. We will show that $U$ is an upper triangular
bidiagonal matrix, i.e., the only entries of $U$ that can be nonzero are the main
diagonal entries $U(i,i)$, for $i = 1:n$, and the upper diagonal entries $U(i,i+1)$,
for $i = 1:(n-1)$.

Following step by step the Cholesky decomposition algorithm from section 6.1.1, we
compute $U(1,1) = \sqrt{A(1,1)}$ and the first row of $U$ as $U(1,k) =
\frac{A(1,k)}{U(1,1)}$, for all $k = 2:n$. Note that $A(1,k) = 0$ if $3 \le k \le n$
since $A$ is a tridiagonal matrix; cf. (6.53). Then, $U(1,k) = 0$ if $3 \le k \le n$,
and therefore the only possible nonzero entries from the first row of $U$ correspond to
an upper triangular bidiagonal form for $U$.

Moreover, since $U(1,2:n) = (U(1,2)\ 0\ \ldots\ 0)$, it follows that
\[
\begin{aligned}
\big(U(1,2:n)\big)\tp U(1,2:n)
&\;=\;
\begin{pmatrix} U(1,2) \\ 0 \\ \vdots \\ 0 \end{pmatrix}
\begin{pmatrix} U(1,2) & 0 & \cdots & 0 \end{pmatrix} \\
&\;=\;
\begin{pmatrix}
(U(1,2))^2 & 0 & \cdots & 0 \\
0 & 0 & \cdots & 0 \\
\vdots & \vdots & & \vdots \\
0 & 0 & \cdots & 0
\end{pmatrix}.
\end{aligned}
\]

Thus, updating the $(n-1) \times (n-1)$ lower right part of $A$ as in (6.28), i.e.,
\[
A(2:n,2:n) \;=\; A(2:n,2:n) \;-\; \big(U(1,2:n)\big)\tp U(1,2:n),
\]
only involves changing the value of $A(2,2)$ to
\[
A(2,2) \;=\; A(2,2) \;-\; (U(1,2))^2,
\]
which preserves the tridiagonal structure of the matrix $A(2:n,2:n)$.

The tridiagonal structure of the updated part of $A$ and the bidiagonal structure of
$U$ will be further preserved as every row of $U$ is computed.

For example, assuming that the updated form $A(i:n,i:n)$ of the matrix $A$ is
tridiagonal after $i-1$ rows of $U$ are computed, the $i$-th row of $U$ is computed as
follows: $U(i,i) = \sqrt{A(i,i)}$ and $U(i,k) = \frac{A(i,k)}{U(i,i)}$, for all $k =
(i+1):n$. Since $A(i:n,i:n)$ is tridiagonal, it follows that $A(i,k) = 0$ if $i+2 \le k
\le n$, and therefore
\[
U(i,k) \;=\; 0, \qquad \forall\, i+2 \le k \le n.
\]
Thus, the only possible nonzero entries from the $i$-th row of $U$ correspond to an
upper triangular bidiagonal form for $U$.

Moreover, since $U(i,i+1:n) = (U(i,i+1)\ 0\ \ldots\ 0)$, we find that
\[
\begin{aligned}
\big(U(i,i+1:n)\big)\tp U(i,i+1:n)
&\;=\;
\begin{pmatrix} U(i,i+1) \\ 0 \\ \vdots \\ 0 \end{pmatrix}
\begin{pmatrix} U(i,i+1) & 0 & \cdots & 0 \end{pmatrix} \\
&\;=\;
\begin{pmatrix}
(U(i,i+1))^2 & 0 & \cdots & 0 \\
0 & 0 & \cdots & 0 \\
\vdots & \vdots & \ddots & \vdots \\
0 & 0 & \cdots & 0
\end{pmatrix}.
\end{aligned}
\]

Thus, updating the $(n-i) \times (n-i)$ lower right part of $A$ as follows as in (6.29),
i.e.,
\[
A(i+1:n,i+1:n) \;=\; A(i+1:n,i+1:n) \;-\; \big(U(i,i+1:n)\big)\tp U(i,i+1:n),
\]
involves only changing the value of $A(i+1,i+1)$ to
\[
A(i+1,i+1) \;=\; A(i+1,i+1) \;-\; (U(i,i+1))^2,
\]
which preserves the tridiagonal structure of the updated $(n-i-1) \times (n-i-1)$
matrix $A(i+1:n,i+1:n)$.

The Cholesky decomposition routine from Table 6.1 simplifies to the pseudocode from
Table 6.4 for tridiagonal symmetric positive definite matrices.

The operation count for the tridiagonal Cholesky decomposition from Table 6.4 is $4n-3$
operations: in each step of the ``for'' loop

% NOTE: the source itself prints "for i = 1 : (n-1)" TWICE in a row here,
% immediately followed by the loop body -- this is an apparent typo in the
% original, reproduced exactly as printed (not corrected).
\begin{center}
\fbox{\begin{minipage}{0.8\linewidth}
for $i = 1 : (n-1)$\\
for $i = 1 : (n-1)$\\
\hspace*{1.5em}$U(i,i) = \sqrt{A(i,i)};\ U(i,i+1) = A(i,i+1)/U(i,i);$\\
\hspace*{1.5em}$A(i+1,i+1) = A(i+1,i+1) - U(i,i+1)^2;$\\
end
\end{minipage}}
\end{center}

we perform 4 operations, plus one more operation for computing $\sqrt{A(n,n)}$, for a
total number of operations of
\begin{equation}
4(n-1) + 1 \;=\; 4n - 3.
\end{equation}
